\section{Metodologia}

A Figura~\ref{fig:pipeline} sintetiza o desenho experimental. A partir de uma única versão auditada do TACO, mantemos a partição e as caixas constantes, derivamos três taxonomias, treinamos cada uma sob protocolo idêntico e aplicamos cinco análises complementares no mesmo conjunto de teste.

\begin{figure}[H]
  \centering
  \includegraphics[width=.94\linewidth]{methodology_pipeline.pdf}
  \caption{Pipeline experimental. Somente a taxonomia varia entre os três ramos de treinamento; imagens, caixas, partição, arquitetura e hiperparâmetros permanecem fixos.}
  \label{fig:pipeline}
\end{figure}

\subsection{Dados, auditoria e partição}

Usamos exclusivamente \texttt{data/annotations.json} do repositório oficial TACO, em formato COCO \cite{lin2014coco}, e somente \textit{bounding boxes}. A auditoria programática confirmou 1.500 imagens, 4.784 objetos, 60 categorias, nenhuma imagem vazia e nenhum arquivo ausente. Sete caixas excediam a borda em no máximo 1,32 pixel; elas foram recortadas apenas na conversão, mantendo o JSON oficial intacto e registrando cada alteração.

A partição foi sorteada uma vez por imagem com semente 42: treino, validação e teste receberam, respectivamente, 1.050/225/225 imagens e 3.359/788/637 objetos. As três taxonomias reutilizam exatamente essa partição. As coordenadas COCO $(x,y,w,h)$ foram convertidas para centro, largura e altura normalizados do YOLO; scripts verificaram limites, dimensões positivas, IDs e igualdade das partições. Imagens com orientação EXIF foram transpostas fisicamente antes da conversão.

\begin{figure}[H]
  \centering
  \includegraphics[width=.94\linewidth]{class_rank_distribution.pdf}
  \caption{Instâncias de treino das 60 classes Fine, ordenadas por frequência. As linhas marcam os limiares da análise de cauda longa.}
  \label{fig:distribution}
\end{figure}

\subsection{Taxonomias}

Fine preserva as 60 categorias oficiais. Material remapeia todas elas para Plastic (24 categorias de origem), Metal (9), Glass (4), Paper/Cardboard (13), Organic (1) e Other (9). Compostos ou itens sem associação natural, como bateria, sapato e embalagens multicamadas, foram enviados a Other; cada decisão e sua justificativa estão em \path{configs/material_mapping.json}. Binary atribui \textit{Litter} a todas as instâncias. As caixas não mudam entre versões. A Figura~\ref{fig:distribution} evidencia a cauda longa que motiva os agrupamentos.

\subsection{Modelo e protocolo experimental}

Treinamos YOLO26n inicializado com pesos COCO, entrada $640\times640$, lote 24, até 60 épocas, \textit{early stopping} com paciência 15, AdamW \cite{loshchilov2019adamw} e taxa inicial $10^{-3}$. Foram mantidas as transformações padrão da implementação Ultralytics, precisão mista e determinismo. Todos os experimentos usaram Python 3.12.10, PyTorch 2.8.0+cu128 e Ultralytics 8.4.157 em uma RTX 3060 Laptop de 6~GB. As sementes 42, 123 e 2026 foram aplicadas às três taxonomias; não houve ajuste específico por configuração.

\subsection{Avaliação e incerteza}

No teste fixo, reportamos precisão, revocação, $F_1$, AP$_{50}$ e AP$_{50:95}$ da implementação Ultralytics; esta última média usa IoU de 0,50 a 0,95 em passos de 0,05. A Tabela principal apresenta média e desvio-padrão entre sementes.

Para decompor o erro, um avaliador controlado faz casamento guloso um-para-um em ordem decrescente de confiança e AP interpolado em 101 pontos. Para cada modelo Fine, preservamos caixas e escores, mas remapeamos predição e verdade-terreno para Fine, Material ou Binary. Como essa implementação não é numericamente idêntica ao avaliador nativo, seus valores só são comparados entre si.

Aplicamos bootstrap percentil pareado \cite{efron1993bootstrap} para obter ICs de 95\%. Fizemos 2.000 reamostragens das 225 imagens de teste, com reposição e os mesmos índices nos dois lados. O bootstrap usa as predições da semente 42; as outras análises agregam três sementes. Para cauda longa, a frequência é calculada apenas no treino: rara ($n\leq19$), média ($20\leq n\leq100$) e frequente ($n>100$).
